% ==== AX ====
\yoschapterstar{Appendix to Chapter X. The Nuclear Space of A. Grothendieck}
\label{chap:AX}

The nuclear operator defined in Chapter~\ref{chap:X}, \ref{sec:X.2} may be extended to locally convex spaces as follows. 

\begin{proposition}\label{prop:AX.0.Proposition.1} Let $X$ be a locally convex linear topological space, and $Y$ a $B$-space. Suppose that there exist an equi-continuous sequence $\{f_{j}^{\prime}\}$ of continuous linear functionals on $X$, a bounded sequence $\{y_{j}\}$ of elements $\in Y$ and a sequence of non-negative numbers $\{c_{j}\}$ with $\sum\limits_{j=1}^{n} c_{j} < \infty$. Then 
\begin{equation}\label{eq:AX.0.1}
T \cdot x = \slim _{n \rightarrow \infty} \sum\limits_{j=1} ^{n} c _{j} \langle x, f _{j} ^{\prime} \rangle y _{j}
\end{equation}
defines a continuous linear operator on $X$ into $Y$. 
\end{proposition}

\begin{proof}
 By the equi-continuity of $\{f_j'\}$, there exists a continuous seminorm $p$ on $X$ such that $\sup |\langle x, f_j' \rangle| \leqq p(x)$ for $x \in X$. Hence, for $m > n$, 
\[
\left\| \sum\limits_{j=n} ^{m} c _{j} \langle x, f _{j} ^{\prime} \rangle y _{j} \right\| \leqq p (x) \sup\limits _{j \geqq 1} \| y _{j} \| \cdot \sum\limits_{j=n} ^{m} c _{j}. \]

This proves that the right hand side of \eqref{eq:AX.0.1} exists and defines a continuous linear operator $T$ on $X$ into the $B$-space $Y$. 
\end{proof}

\begin{definition}\label{def:AX.0.Definition.1} An operator $T$ of the form \eqref{eq:AX.0.1} is said to be a nuclear operator on $X$ into the $B$-space $Y$. 
\end{definition}

\begin{corollaryu} $A$ nuclear operator $T$ is a compact operator in the sense that it maps a neighbourhood of $O$ of $X$ into a relatively compact set of $Y$. 

\end{corollaryu}

\begin{proof}
 We define 
\[
T _{n} \cdot x = \sum\limits_{j=1} ^{n} c _{j} \langle x, f _{j} ^{\prime} \rangle y _{j}.
\]
$T_{n}$ is compact, since the image by $T_{n}$ of the set $V = \{x; p(x) \leqq 1\}$ of $X$ is relatively compact in $Y$. On the other hand, we have 
\[
\left\| T x - T _{n} x \right\| = \left\| \sum\limits_{j=n+1} ^{\infty} c _{j} \langle x, f _{j} ^{\prime} \rangle y _{j} \right\| \leqq p (x) \sup\limits _{j \geqq 1} \left\| y _{j} \right\| \sum\limits_{j=n+1} ^{\infty} c _{j},
\]
and so $T_{n}x$ converges to $Tx$ strongly and uniformly on $V$. Hence the operator $T$ is compact. 
\end{proof}

As was proved in Chapter~\ref{chap:X}, \ref{sec:X.2}, we have a typical example of a nuclear operator: 

\begin{examplen} Let $K$ be a compact subset of $\mathbb{R}$. Then, for $(k-j)>n$, the identity mapping $T$ of $H_{0}^{k}(K)$ into $H_{0}^{j}(K)$ is a nuclear operator.

\end{examplen}
\begin{proposition}\label{prop:AX.0.Proposition.2} Let $X$ be a locally convex linear topological space, and $V$ a convex balanced neighbourhood of $0$ of $X$. Let $p_{V}(x) = \inf\limits_{x/\lambda \in V, \lambda > 0} \lambda$ be the Minkowski functional of $V$. $p_{V}$ is a continuous semi-norm on $X$. Set 
\[
N _{V} = \{x \in X; p _{V} (x) = 0 \} = \{x \in X; \lambda x \in V \text{ for all } \lambda > 0 \}.
\]
\end{proposition}

Then $N_V$ is a closed linear subspace of $X$, and the quotient space $X_V = X / N_V$ is a normed linear space by the norm 
\begin{equation}\left\| \tilde {x} \right\| _{V} = p _{V} (x), \text{ where } \tilde {x} \text{ is the residue class mod } N _{V} \text{ containing the element } x. \label{eq:AX.0.2}\end{equation}

\begin{proof}
 Let $(x - x_1) \in N_V$. Then $p_V(x_1) \leqq p_V(x) + p_V(x_1 - x) = p_V(x)$, and similarly $p_V(x) \leqq p_V(x_1)$. Thus $p_V(x) = p_V(x_1)$ if $x$ and $x_1$ are in the same residue class mod $N_V$. We have $\|\tilde{x}\|_V \geqq 0$ and $\|\tilde{0}\|_V = 0$. If $\|\tilde{x}\|_V = 0$, then $x \in \tilde{x}$ implies $x \in N_V$ and so $\tilde{x} = \tilde{0}$. The triangle inequality is proved by $\|\tilde{x} + \tilde{y}\|_V = p_V(x + y) \leqq p_V(x) + p_V(y) = \|\tilde{x}\|_V + \|\tilde{y}\|_V$. We have also $\|\alpha \tilde{x}\|_V = p_V(\alpha x) = |\alpha | p_V(x) = |\alpha \|\tilde{x}\|_V$. 
\end{proof}

\begin{corollaryu} By the equivalence 
\begin{equation}\label{eq:AX.0.3}
\left(p _ {V _{1}} \leqq p _ {V _{2}}\right) \leftrightarrow \left(V _{2} \subseteq V _{1}\right),
\end{equation}
we can define the canonical mapping 
\[
X _ {V _{2}} \rightarrow X _ {V _{1}} (\text{ when } V _{2} \subseteq V _{1})
\]
by associating the residue class $\tilde{x}_{V_{1}}\pmod{N_{V_{1}}}$ containing $x$ to the residue class $\tilde{x}_{V_{1}}\pmod{N_{V_{1}}}$ containing $\tilde{x}$. The mapping thus obtained is continuous, since 
\[
\left\| \tilde {x} _ {V _{1}} \right\| _ {V _{1}} = p _ {V _{1}} (x) \geqq p _ {V _{1}} (x) = \left\| \tilde {x} _ {V _{1}} \right\| _ {V _{1}}.
\]
\end{corollaryu}

We are now ready to give the notion of a nuclear space, introduced in analysis by A. \textsc{Grothendieck} \cite{grothendieck88}. 

\begin{definition}\label{def:AX.0.Definition.2} $A$ locally convex linear topological space $X$ is said to be a nuclear space, if, for any convex balanced neighbourhood $V$ of $0$, there exists another convex balanced neighbourhood $U \subseteq V$ of $0$ such that the canonical mapping 
\begin{equation}\label{eq:AX.0.4}
T \colon X _{U} \to \hat {X} _{V}
\end{equation}
is nuclear. Here $\hat{X}_{V}$ is the completion of the normed linear space $X_{V}$. 
\end{definition}

\begin{example}\label{ex:AX.0.Example.1} Let $\mathbb{R}^{A}$ be the topological product of real number field $\mathbb{R}$ in such a way that $\mathbb{R}^{A}$ is the totality of real-valued finite functions $x(a)$ defined on $A$ and topologized by the system of semi-norms 
\begin{equation}\label{eq:AX.0.5}
p_{a} (x) = | x (a) |, a \in A.
\end{equation}
\end{example}

Then $\mathbb{R}^{A}$ is a nuclear space. 

\begin{proof}
 $N_{V}$ is the totality of functions $x(a) \in \mathbb{R}^{A}$ such that, for some finite set $\{a_{j} \in A; j = 1, 2, \ldots, n\}$, $x(a_{j}) = 0$ ( $j = 1, 2, \ldots, n$). Hence $X_{V} = \mathbb{R}^{A}/N_{V}$ is equivalent to the space of functions $x_{V}(a)$ such that $x_{V}(a) = 0$ for $a \neq a_{j}$ ( $j = 1, 2, \ldots, n$) and normed by 
\[
\left\| x _{V} (a) \right\| _{V} = \sup\limits _{1 \leqq j \leqq n} | x (a _{j}) |. \]

We take $N_U$ the totality of functions $x(a) \in \mathbb{R}^{A}$ such that $x(a_\alpha) = 0$ for $\alpha \in A'$ where $A'$ is any finite set of integers containing $1, 2, \ldots, n$. Thus, for $U \subseteq V$, the canonical mapping $X_U = \mathbb{R}^{A} / N_U \to \mathbb{R}^{A} / N_V = X_V$ is nuclear. For the mapping is a continuous linear mapping with a finite-dimensional range. 
\end{proof}

\begin{example}\label{ex:AX.0.Example.2} $A$ nuclear $B$-space $X$ must be of finite dimension. 
\end{example}

\begin{proof}
 Since $X = X_{V}$ for any convex balanced neighbourhood $V$ of $0$ of a $B$-space, the compactness of the identity mapping $X \to X$ implies that $X$ is of finite dimension by F. \textsc{Riesz}' theorem in Chapter~\ref{chap:III}, \ref{sec:III.2}. 
\end{proof}

\begin{example}\label{ex:AX.0.Example.3} Let $K$ be a compact subset of $\mathbb{R}^{n}$. Then the space $\mathfrak{D}_K(\mathbb{R}^{n})$ introduced in Chapter~\ref{chap:I}, \ref{sec:I.1} is a nuclear space. 
\end{example}

\begin{proof}
 As in Chapter~\ref{chap:I}, \ref{sec:I.1}, let 
\[
\phi_{K,k} (f) = \sup\limits _{x \in K, | s | \leqq k} | D ^{s} f (x) |
\]
be one of the semi-norms which define the topology of $\mathfrak{D}_K(\mathbb{R}^{n})$. Let $V_k = \{f \in \mathfrak{D}_K(\mathbb{R}^{n}); p_{K,k}(f) \leqq 1\}$. Then $N_{V_{k}}$ is $\{0\}$, and $X_{V_{k}} = X / N_{V_{k}} = \mathfrak{D}_K(\mathbb{R}^{n}) / N_{V_{k}}$ is precisely the space $\mathfrak{D}_K(\mathbb{R}^{n})$ normed by $p_{K,k}$. If $(k-j) > n$, then it is easy to prove, as in the example following the Corollary of Definition~\ref{def:AX.0.Definition.1} above, that the canonical mapping of $X_{V_{k}}$ into $X_{V_{j}}$ is a nuclear transformation. Hence $\mathfrak{D}_K(\mathbb{R}^{n})$ is a nuclear space. 
\end{proof}

\begin{theorem}\label{thm:AX.0.Theorem.1} $A$ locally convex linear topological space $X$ is nuclear, iff, for any convex balanced neighbourhood $V$ of $0$, the canonical mapping $X \rightarrow \hat{X}_{V}$ is nuclear. 
\end{theorem}

\begin{proof}
 Necessity. Let $U \subseteq V$ be a convex balanced neighbourhood of $0$ of $X$ such that the canonical mapping $X_{U} \to \hat{X}_{V}$ is nuclear. The canonical mapping $T: X \to \hat{X}_{V}$ is the product of the canonical mapping $X \to \hat{X}_{U}$ and the canonical nuclear transformation $X_{U} \to \hat{X}_{V}$. Hence $T$ must be a nuclear transformation. 

Sufficiency. Let the canonical mapping $T: X \to \hat{X}_{V}$ be given by a nuclear transformation 
\[
T x = \sum\limits_{j=1} ^{\infty} c _{j} \langle x, f _{j} ^{\prime} \rangle y _{j}.
\]

For any $\alpha > 0$, the set $\{x \in X; |\langle x, f_j' \rangle| \leqq \alpha$ for $j = 1, 2, \ldots\}$ is a convex balanced neighbourhood $U_{\alpha}$ of $0$ of $X$, because of the equi-continuity of $\{f_j'\} \subseteq X'$. Moreover, 
$\| Tx\| _V = \left\| \sum\limits_{j}c_j\langle x,f_j'\rangle y_j\right\| _V\leqq \alpha \sup\limits_{j}\| y_j\| _V\sum\limits_{j}c_j$ whenever $x\in U_{\alpha}$ 

Let $\alpha$ be so small that the right hand side is $< 1$. Then $\|Tx\|_V < 1$ and $U_{\alpha} \subseteq V$. Each $f_j'$ may be considered as belonging to the dual space $X_{U_{\alpha}}'$, and so 
\[
T x = T z = \sum\limits_{j} c _{j} \left\langle x, f _{j} ^{\prime} \right\rangle y _{j} \quad \text{ whenever } \quad (x - z) \in N _ {U _{\alpha}}.
\]

Thus the canonical mapping $X_{U_{\alpha}} \to \hat{X}_V$ is given by a nuclear transformation 
\[
\tilde {x} _ {U _{\alpha}} \rightarrow \sum\limits_{j} c _{j} \left\langle \tilde {x} _ {U _{\alpha}}, f _{j} ^{\prime} \right\rangle y _{j}.
\]
\end{proof}

\begin{theorem}\label{thm:AX.0.Theorem.2} Let a locally convex linear topological space $X$ be nuclear. Then, for any convex balanced neighbourhood $V$ of $0$ of $X$, there exists a convex balanced neighbourhood $W \subseteq V$ of $0$ of $X$ such that $\hat{X}_{W}$ is a Hilbert space. 
\end{theorem}

\begin{proof}
 The nuclear canonical mapping $X_{U} \to \hat{X}_{V} (U \subseteq V)$ defined by 
\[
T \tilde {x} _{U} = \sum\limits_{j} c _{j} \left\langle \tilde {x} _{U}, f _{j} ^{\prime} \right\rangle y _{j}
\]
is factored as the product of the two mappings 
$\alpha : X_U \to (l^2)$ and $\beta : (l^2) \to \hat{X}_V$, where $\alpha$ is given by 
\[
\tilde {x} _{U} \rightarrow \left\{c _{j} ^{1/2} \langle \tilde {x} _{U}, f _{j} ^{\prime} \rangle \right\} \quad \text{ and } \quad \beta \text{ is given by } \{\xi_{j} \} \rightarrow \sum\limits_{j} c _{j} ^{1/2} \xi_{j} y _{j}. \]

The continuity of $\alpha$ is clear from 
\[
\sum\limits_{j} \left| c _{j} ^{1/2} \langle \tilde {x} _{U}, f _{j} ^{\prime} \rangle \right| ^{2} \leqq \left(\sup\limits _{j} \left\| f _{j} ^{\prime} \right\| \cdot \left\| \tilde {x} _{U} \right\| _{U}\right) ^{2} \cdot \sum\limits_{j} c _{j},
\]
and that of $\beta$ is proved by 
\[
\left\| \sum\limits_{j} c _{j} ^{1/2} \xi_{j} y _{j} \right\| _{V} ^{2} \leqq \sum\limits_{j} c _{j} \| y _{j} \| _{V} ^{2} \cdot \sum\limits_{j} | \xi_{j} | ^{2} \leqq \sup\limits _{j} \| y _{j} \| _{V} ^{2} \cdot \| \{\xi_{j} \} \| _ {l ^{n}} ^{2} \cdot \sum\limits_{j} c _{j}.
\]

Let $U_{2}$ be the inverse image in $(l^{2})$ by $\beta$ of the unit sphere of $\hat{X}_{V}$. Then $U_{2}$ is a neighbourhood of $0$ of $(l^{2})$ and so contains a sphere $S$ of centre $0$ of $(l^{2})$. 

Let $W$ be the inverse image in $X$ of $S$ by the continuous mapping $\tilde{\alpha}$ defined as the product of continuous canonical mapping $X \to X_U$ and the continuous mapping $\alpha: X_U \to (l^2)$. Then clearly $W \subseteq V$ and, for any $\tilde{x}_W \in X_W$, 
\[
\left\| \tilde {x} _{W} \right\| _{W} = \inf\limits _{x / \lambda \in W, \lambda > 0} \lambda = \inf\limits _ {\tilde {\alpha} x / \lambda \in S, \lambda > 0} \lambda = \left\| \tilde {\alpha} x \right\| _ {L ^{2}} (\text{ the radius of } S).
\]

Since $\|\|_{l^2}$ is the norm in the Hilbert space $(l^2)$, $X_W$ is a pre-Hilbert space. 
\end{proof}

\begin{corollaryu} Let $X$ be a locally convex nuclear space. Then, for any convex balanced neighbourhood $V$ of $0$ of $X$, there exist convex balanced neighbourhoods $W_{1}$ and $W_{2}$ of $0$ of $X$ with the properties: 

$W_{2} \subseteq W_{1} \subseteq V, \hat{X}_{W_{1}}$ and $\hat{X}_{W_{1}}$ are Hilbert spaces and the canonical 
mappings $X \to \hat{X}_{W_2}, \hat{X}_{W_2} \to \hat{X}_{W_1}, \hat{X}_{W_1} \to \hat{X}_V$ are all nuclear. 
\end{corollaryu}

Therefore, a nuclear space $X$ has a fundamental system $\{V_{\alpha}\}$ of neighbourhoods of $0$ such that the spaces $\hat{X}_{V_{\alpha}}$ are Hilbert spaces. 

Further Properties of Nuclear Spaces. It can be proved that: 

$1$. A linear subspace and a factor space of a nuclear space are also nuclear. 

$2$. The topological vector product of a family of nuclear spaces and the inductive limit of a sequence of nuclear spaces are also nuclear. 

$3$. The strong dual of the inductive limit of a sequence of nuclear spaces, each of which is an $F$-space, is also nuclear. 

For the proof, see the book by \textsc{Grothendieck} \cite{grothendieck87} referred to above, p. $47$. As a consequence of $2$., the space $\mathfrak{D}(\mathbb{R}^{n})$, which is the inductive limit of the sequence $\{\mathfrak{D}_{K_{r}}(\mathbb{R}^{n}); r = 1, 2, \ldots\}$ (here $K_r$ is the sphere $|x| \leqq r$ of $\mathbb{R}^{n}$), is nuclear. Hence, by $3$., the space $\mathfrak{D}(\mathbb{R}^{n})'$ is also nuclear. The spaces $\mathfrak{E}(\mathbb{R}^{n}), \mathfrak{E}(\mathbb{R}^{n})'$, $\mathfrak{S}(\mathbb{R}^{n})$ and $\mathfrak{S}(\mathbb{R}^{n})'$ are also nuclear spaces. 

The importance of the notion of the nuclear space has recently been stressed by R. A. \textsc{Minlos} \cite{minlos202}. He has proved the following generalization of Kolmogorov's extension theorem of measures: 

Let $X$ be a nuclear space whose topology is defined through a countable system of convex balanced neighbourhoods of $0$. Let $X'$ be the strong dual space of $X$. A cylinder set of $X'$ is defined as a set of the form 
\[
Z ^{\prime} = \{f ^{\prime} \in X ^{\prime}; a _{i} < \langle x _{i}, f ^{\prime} \rangle < b _{i} (i = 1, 2, \dots , n) \}.
\]

Suppose there is given a set function $\mu_{0}$, defined and $\geqq 0$ for all cylinder sets. Let $\mu_{0}$ be $\sigma$-additive for those cylinder sets $Z'$ with fixed $x_{1}, x_{2}, \ldots, x_{n}$. Then, under a compatibility condition and a continuity condition, there exists a uniquely determined extension of $\mu_{0}$ which is $\sigma$-additive and $\geqq 0$ for all sets of the smallest $\sigma$-additive family of sets of $X'$ containing all the cylinder sets of $X'$. 

For a detailed proof and applications of this result, see I. M. \textsc{Gelfand}-N. Y. \textsc{Vilenkin} \cite{gelfand77}.\yosstarfootnote{See also Supplementary Notes, p.~\pageref{notes:AX}.} 
